\documentclass[../main.tex]{subfiles}
\begin{document}
%=========================================================
% ĐÁP ÁN VÀ HƯỚNG DẪN GIẢI - DÀNH CHO GIÁO VIÊN
% Phiếu bài tập Toán tư duy cấp độ A - HK1 Tuần 17 (Ôn tập phép nhân, ôn tập bài toán có lời văn)
%=========================================================
\providecommand{\DAbai}{}\renewcommand{\DAbai}[1]{\par\smallskip\noindent{\color{mauchinh}\bfseries Bài #1.}\ }
\providecommand{\DAphan}{}\renewcommand{\DAphan}[1]{\par\medskip\noindent{\bfseries\color{orange!80!black}#1}\par\nopagebreak}
% Bài giải mẫu theo đúng cách trình bày lớp 2: \DAgiai{câu lời giải}{phép tính}{đáp số}
\providecommand{\DAgiai}{}\renewcommand{\DAgiai}[3]{\par\noindent\hspace*{0.6cm}\begin{minipage}[t]{15.5cm}\textit{Bài giải}\\ #1\\ \hspace*{0.8cm}#2\\ Đáp số: #3\end{minipage}\par}
% Bài nhiều bước: \DAnhieu{\DAb{câu lời giải}{phép tính}\DAb{câu lời giải}{phép tính}}{đáp số}
\providecommand{\DAb}{}\renewcommand{\DAb}[2]{#1\\ \hspace*{0.8cm}#2\\ }
\providecommand{\DAnhieu}{}\renewcommand{\DAnhieu}[2]{\par\noindent\hspace*{0.6cm}\begin{minipage}[t]{15.5cm}\textit{Bài giải}\\ #1Đáp số: #2\end{minipage}\par}

\par\medskip
{\centering
        {\color{mauchinh!120}\bfseries\Large ĐÁP ÁN VÀ HƯỚNG DẪN GIẢI}\\[3pt]
        {\color{mauchinh!120}\bfseries\Large (DÀNH CHO GIÁO VIÊN) - ÔN TẬP PHÉP NHÂN VÀ BÀI TOÁN CÓ LỜI VĂN}\par}
\vspace{6pt}

\noindent\textit{Lưu ý: (1) Phần A và B chỉ dùng bảng nhân $2$ và bảng nhân $5$, đúng chương trình Toán lớp 2. (2) Các bài phần C mô phỏng \textbf{dạng} đề của các kì thi Kangaroo, SASMO, Olympic Toán Singapore, Violympic, không phải trích nguyên văn đề thi; một số bài cần thử nhiều trường hợp hoặc cộng các số hạng bằng nhau (HS chưa học bảng nhân $3$, $4$, $6$). (3) Bài giải theo mẫu: câu lời giải, phép tính kèm đơn vị trong ngoặc, đáp số; số dòng kẻ trên phiếu tương ứng với số bước giải.}

%%%%%%%%%%%%%%%%%%%%%%%%%%%%%%%%%%%%%%%%%%%%%%%%%%%%%%%%%%
\section*{BUỔI 1. ÔN TẬP PHÉP NHÂN}
\DAphan{Phần A}
\DAbai{1} a) $2+2+2+2+2=2\times5=10$; b) $5+5+5+5=5\times4=20$; c) $2\times6=2+2+2+2+2+2=12$; d) $5\times4=5+5+5+5=20$.

\DAbai{2} a) $2$ và $6$ gọi là \textbf{thừa số}; $12$ gọi là \textbf{tích}. b) Từ trái sang phải: cột 1 tích là $14$; cột 2 thừa số còn thiếu là $6$ ($5\times6=30$); cột 3 thừa số còn thiếu là $2$ ($2\times9=18$); cột 4 tích là $50$; cột 5 thừa số còn thiếu là $8$ ($2\times8=16$).

\DAbai{3} Hàng $2\times$ số: $6;\ 10;\ 14;\ 18$ (các ô theo thứ tự $3,5,7,9$). Hàng $5\times$ số: $15;\ 25;\ 35;\ 45$.

\DAbai{4} $2\times7\to14$; $5\times6\to30$; $2\times9\to18$; $5\times8\to40$; $2\times5\to10$.

\DAbai{5} $12>10$; $15<16$; $20=20$; $35>18$; $2\times5=5\times2$ (đổi chỗ các thừa số, tích không đổi); $30=30$ (vì $5\times5+5=25+5=30$).

\DAbai{6} Hình 1: $4$ nhóm, mỗi nhóm $5$ chấm: $5+5+5+5=20$; $5\times4=20$. Hình 2: $5$ nhóm, mỗi nhóm $2$ chấm: $2+2+2+2+2=10$; $2\times5=10$. Hình 3: $3$ nhóm, mỗi nhóm $5$ chấm: $5+5+5=15$; $5\times3=15$.

\DAphan{Phần B}
\DAbai{7} $7;\ 7;\ 9;\ 9;\ 5;\ 10$.

\DAbai{8} a) $30+14=44$; b) $16-9=7$; c) $18+25=43$; d) $35-18=17$; e) $50-23=27$; f) $20+37=57$. \textit{Lưu ý:} cả sáu biểu thức đều có dạng ``tích $\pm$ một số'', HS tính tích trước (đúng thứ tự từ trái sang phải).

\DAbai{9} a) $8,\ 10,\ 14$ (đếm thêm $2$); b) $35,\ 30,\ 20$ (đếm bớt $5$); c) $20,\ 24$ (đếm thêm $4$); d) $22,\ 27$ (đếm thêm $5$); e) $16,\ 32$ (mỗi số bằng $2$ lần số liền trước); f) $16,\ 22$ (số cộng thêm tăng dần $1,2,3,4,5,6$).

\DAbai{10} Số vào $3\to13$; $6\to28$; $7\to33$. Số ra $23\to$ số vào $5$ (vì $5\times5-2=23$); số ra $48\to$ số vào $10$; số ra $3\to$ số vào $1$. \textit{Gợi ý:} với câu hỏi ngược, HS cộng $2$ rồi tìm số nhân với $5$ ra kết quả đó.

\DAbai{11} a) Đ; b) S, sửa $5\times8=40$; c) Đ ($30-10=20$); d) S, sửa $2\times9+5=18+5=23$.

\DAphan{Phần C - Thử thách}
\DAbai{12} a) $2$ và $5$; b) $4$ và $5$; c) $3$ và $4$. \textit{Cách làm:} liệt kê các cặp số có tích đã cho (ví dụ $20=2\times10=4\times5$) rồi kiểm tra tổng.

\DAbai{13} a) $\bigstar=5$ (vì $5+5+5+5=20$), $\heartsuit=7$ (vì $5\times7=35$), $\bigstar+\heartsuit=12$. b) $\clubsuit=6$, $\spadesuit=5$ (vì $6\times5=30$), $\clubsuit+\spadesuit=11$.

\DAbai{14} a) $10,\ 20,\ 30,\ 40,\ 50$ (các số có chữ số hàng đơn vị là $0$). b) Số bí mật là $30$.

\DAbai{15} a) $2\times5+5\times5=10+25=35$. b) Ghép cặp: $(1+10)+(2+9)+(3+8)+(4+7)+(5+6)=11+11+11+11+11=55$. c) $(2+12)+(4+10)+(6+8)=14+14+14=42$.

\DAbai{16} a) $4+4+4=12$ que. b) Hình vuông đầu cần $4$ que, mỗi hình vuông tiếp theo chỉ cần thêm $3$ que: $4+3+3=10$ que. c) $4+3+3+3+3=16$ que.

\DAbai{17} $9$ hình vuông nhỏ ($1\times1$) $+\ 4$ hình vuông ($2\times2$) $+\ 1$ hình vuông ($3\times3$) $=14$ hình vuông.

%%%%%%%%%%%%%%%%%%%%%%%%%%%%%%%%%%%%%%%%%%%%%%%%%%%%%%%%%%
\section*{BUỔI 2. ÔN TẬP BÀI TOÁN CÓ LỜI VĂN}
\DAphan{Phần A}
\DAbai{1} a) $5\times6=30$; đáp số $30$ quả cam. b) $2\times9=18$; đáp số $18$ bánh xe.

\DAbai{2} a) B (Số chiếc bút $7$ bạn mua là:). b) C (Số quả táo còn lại là:).

\DAbai{3} Đề 1 $\to2\times7$; đề 2 $\to15-6$; đề 3 $\to9+5$; đề 4 $\to5\times4$.

\DAbai{4} \DAgiai{Số cây tổ 2 trồng được là:}{$24+9=33$ (cây)}{$33$ cây.}
\DAgiai{Số viên bi An có là:}{$45-17=28$ (viên)}{$28$ viên bi.}

\DAphan{Phần B}
\DAbai{5} Thứ tự đúng: \textbf{D, A, B, C, E}.

\DAbai{6} \DAnhieu{\DAb{Số cái bánh mẹ mua là:}{$5\times4=20$ (cái)}\DAb{Số cái bánh còn lại là:}{$20-7=13$ (cái)}}{$13$ cái bánh.}
\textit{Tóm tắt mẫu:} Mỗi hộp: $5$ cái; $4$ hộp: ? cái; đã ăn: $7$ cái; còn lại: ? cái.

\DAbai{7} a) Bài giải của An \textbf{sai}. b) An sai ở \textbf{phép tính} (đáng lẽ phải nhân vì mỗi bàn có $2$ bạn, có $7$ bàn như thế), kéo theo đáp số sai. c) \DAgiai{Số bạn ngồi là:}{$2\times7=14$ (bạn)}{$14$ bạn.}

\DAbai{8} \DAnhieu{\DAb{Số kẹo của Hà là:}{$5\times3=15$ (cái)}\DAb{Mai có nhiều kẹo hơn Hà số cái là:}{$17-15=2$ (cái)}}{Mai nhiều hơn, $2$ cái kẹo.}

\DAbai{9} (Đáp án mẫu) Đề toán: ``Mỗi đĩa có $5$ cái kẹo. Hỏi $7$ đĩa như thế có tất cả bao nhiêu cái kẹo?'' \DAgiai{Số cái kẹo $7$ đĩa có là:}{$5\times7=35$ (cái)}{$35$ cái kẹo.} Chấp nhận mọi đề đúng ngữ cảnh và đúng phép tính $5\times7$.

\DAbai{10} \DAnhieu{\DAb{Số tiền $3$ chiếc bút là:}{$2\times3=6$ (nghìn đồng)}\DAb{Số tiền Nam phải trả là:}{$6+5=11$ (nghìn đồng)}\DAb{Số tiền cô bán hàng trả lại Nam là:}{$20-11=9$ (nghìn đồng)}}{$9$ nghìn đồng.}

\DAphan{Phần C - Thử thách}
\DAbai{11} Thử: $5$ gà, $2$ thỏ có $5\times2+2\times4=10+8=18$ chân (chưa đủ); $4$ gà, $3$ thỏ có $4\times2+3\times4=8+12=20$ chân (đủ). Đáp số: $4$ con gà và $3$ con thỏ.\par
\noindent\textit{Gợi ý cho HS khá giỏi:} giả sử cả $7$ con đều là gà thì có $14$ chân, còn thiếu $6$ chân; mỗi lần đổi một gà thành một thỏ thì thêm $2$ chân, cần đổi $3$ lần, nên có $3$ thỏ và $4$ gà.

\DAbai{12} Các cây trồng tại các vị trí $0,\ 2,\ 4,\ 6,\ 8,\ 10$ m nên có $6$ cây. Đáp số: $6$ cây. \textit{Lưu ý:} có $5$ khoảng cách nhưng có $6$ cây vì trồng cả hai đầu.

\DAbai{13} Bảng: ngày 1: $3$ m, $2$ m; ngày 2: $5$ m, $4$ m; ngày 3: $7$ m, $6$ m; ngày 4: sáng lên $6+3=9$ m, đã tới (vượt) đỉnh cột $8$ m nên không bị tụt nữa. Đáp số: $4$ ngày. \textit{Lỗi thường gặp:} trả lời $3$ ngày vì quên rằng ở ngày cuối ốc sên lên tới đỉnh thì không bị tụt xuống nữa.

\DAbai{14} Bảng: em $3$ tuổi, anh $7$ tuổi, tổng $10$; em $4$ tuổi, anh $8$ tuổi, tổng $12$ (đủ). Đáp số: anh $8$ tuổi, em $4$ tuổi.

\DAbai{15} Có $6$ đoạn thẳng nên có $6$ cái bắt tay: $AB,\ AC,\ AD,\ BC,\ BD,\ CD$ (hoặc $3+2+1=6$). Đáp số: $6$ cái bắt tay.

\begin{center}
       \faGg\faGg\quad Hết \quad\faGg\faGg
\end{center}
\end{document}
